\chapter{Technologies et implémentation}
\label{chap:techno}

Ce chapitre justifie les choix technologiques selon une grille uniforme
(pourquoi, comment, bénéfices, limites, alternatives) puis détaille les points
saillants de l'implémentation.

\section{Présentation de la pile technologique}
Le tableau~\ref{tab:stack} synthétise la pile par niveau.

\begin{table}[H]
\centering\caption{Pile technologique de DevPilot AI}\label{tab:stack}\small
\begin{tabularx}{\textwidth}{@{}>{\bfseries}l l X@{}}
\toprule
Niveau & Technologie & Rôle \\
\midrule
Backend & FastAPI, Pydantic & API REST et \emph{streaming}, validation \\
 & SQLAlchemy, Alembic & ORM asynchrone et migrations \\
Base de données & PostgreSQL 16 & État relationnel et traces \\
Messagerie & Redis 7 & Courtier Celery et \emph{cache} \\
Asynchrone & Celery & Traitement documentaire \\
Vectoriel & Qdrant 1.12 & Recherche HNSW (cosinus) \\
Modèles IA & Ollama, Gemini & \emph{Embeddings} locaux, génération \\
Frontend & Next.js 15, React 19, TypeScript & Interface, chat, tableaux de bord \\
 & React Query, Tailwind, Framer Motion, React Flow & État serveur, design, animations, graphes \\
Observabilité & Prometheus, Grafana & Métriques et tableaux de bord \\
Déploiement & Docker, Docker Compose & Conteneurisation (8 services) \\
\bottomrule
\end{tabularx}
\end{table}

\section{Technologies backend}

\subsection{FastAPI}
\textbf{Pourquoi.} Framework Python asynchrone hautes performances avec validation
native. \textbf{Comment.} Routeurs versionnés, dépendances d'autorisation,
\texttt{StreamingResponse} pour le SSE. \textbf{Bénéfices.} Concurrence par E/S
non bloquantes, documentation OpenAPI automatique. \textbf{Limites.} Écosystème
plus jeune que des frameworks synchrones établis. \textbf{Alternatives.} Django
REST Framework, Flask~\cite{fastapi}.

\subsection{PostgreSQL}
\textbf{Pourquoi.} SGBD relationnel robuste et transactionnel, idéal pour l'état et
les traces structurées. \textbf{Comment.} Onze tables, clés UUID, support JSON
pour les champs semi-structurés (\emph{previews} d'agents). \textbf{Bénéfices.}
Intégrité, requêtes analytiques, maturité. \textbf{Limites.} Inadapté au stockage
vectoriel natif, d'où le recours à Qdrant. \textbf{Alternatives.} MySQL,
SQLite~\cite{postgres}.

\subsection{SQLAlchemy et Alembic}
\textbf{Pourquoi.} ORM expressif et migrations versionnées. \textbf{Comment.}
Modèles déclaratifs asynchrones (asyncpg), migrations Alembic. \textbf{Bénéfices.}
Typage, portabilité, évolution maîtrisée du schéma. \textbf{Limites.} Courbe
d'apprentissage de l'API asynchrone. \textbf{Alternatives.} Tortoise ORM,
SQLModel.

\subsection{Redis}
\textbf{Pourquoi.} Courtier de messages à faible latence et magasin de résultats.
\textbf{Comment.} \emph{Broker} et \emph{backend} de Celery, socle de \emph{cache}.
\textbf{Bénéfices.} Rapidité, simplicité opérationnelle. \textbf{Limites.}
Persistance mémoire à dimensionner. \textbf{Alternatives.} RabbitMQ,
Amazon SQS~\cite{redis}.

\subsection{Celery}
\textbf{Pourquoi.} Standard éprouvé du traitement de tâches asynchrones en Python.
\textbf{Comment.} Tâche \texttt{process\_document\_task} (extraction, découpage,
vectorisation, indexation), relancée jusqu'à trois fois. \textbf{Bénéfices.}
Découplage du chemin critique, montée en charge des \emph{workers}.
\textbf{Limites.} Complexité opérationnelle accrue. \textbf{Alternatives.} RQ,
Dramatiq~\cite{celery}.

\subsection{Qdrant}
\textbf{Pourquoi.} Base vectorielle spécialisée et auto-hébergeable.
\textbf{Comment.} Collection indexée en HNSW, distance cosinus, vecteurs de
dimension~768. \textbf{Bénéfices.} Recherche par plus proches voisins en temps
quasi logarithmique~\cite{malkov2018}. \textbf{Limites.} Composant additionnel à
exploiter. \textbf{Alternatives.} Milvus, Weaviate, pgvector~\cite{qdrant}.

\subsection{Ollama}
\textbf{Pourquoi.} Exécution locale de modèles d'\emph{embedding}, fonctionnement
hors-ligne. \textbf{Comment.} Modèle \texttt{nomic-embed-text}, vectorisation par
lots. \textbf{Bénéfices.} Confidentialité, absence de coût d'API pour les
\emph{embeddings}. \textbf{Limites.} Dépendance aux ressources de l'hôte.
\textbf{Alternatives.} \emph{embeddings} OpenAI, Sentence-Transformers~\cite{ollama,reimers2019}.

\subsection{Gemini}
\textbf{Pourquoi.} Génération de haute qualité avec \emph{streaming} natif.
\textbf{Comment.} \texttt{gemini-2.5-flash} (repli \texttt{gemini-2.0-flash}),
température 0,2, 1000 jetons maximum, flux SSE via \texttt{streamGenerateContent}.
\textbf{Bénéfices.} Fluidité, faible latence, repli automatique.
\textbf{Limites.} Dépendance à un service distant. \textbf{Alternatives.}
OpenAI GPT, modèles locaux via Ollama~\cite{gemini}.

\section{Technologies frontend}

\subsection{Next.js}
\textbf{Pourquoi.} \emph{App Router} React 19, rendu hybride, excellente
expérience développeur. \textbf{Comment.} Pages par route, modules fonctionnels,
chargement différé des visualisations lourdes. \textbf{Bénéfices.} Découpage
automatique des \emph{bundles}, \emph{streaming} côté client. \textbf{Limites.}
Couplage fort à l'écosystème React. \textbf{Alternatives.} Remix,
SvelteKit~\cite{nextjs}.

\subsection{TypeScript}
\textbf{Pourquoi.} Sûreté de typage à grande échelle. \textbf{Comment.} Mode
strict, zéro \texttt{any} sur l'ensemble du code applicatif. \textbf{Bénéfices.}
Détection précoce des erreurs, refactorisations sûres. \textbf{Limites.} Surcoût
de typage initial. \textbf{Alternatives.} JavaScript + JSDoc, Flow.

\subsection{React Query}
\textbf{Pourquoi.} Gestion robuste de l'état serveur. \textbf{Comment.} Requêtes
mises en cache, dédupliquées et invalidées (inspecteur de traces, observatoire,
tableau de bord personnel). \textbf{Bénéfices.} Suppression des conditions de
course, standardisation des états de chargement et d'erreur. \textbf{Limites.}
Migration progressive depuis les appels manuels. \textbf{Alternatives.} SWR,
RTK Query.

\subsection{Tailwind CSS et système de design « Horizon »}
\textbf{Pourquoi.} Cohérence visuelle et productivité. \textbf{Comment.} Jetons CSS
(thème clair/sombre) mappés sur des utilitaires sémantiques ; primitives
réutilisables (Button, Card, DataCard, MetricCard, Badge, Dialog). \textbf{Bénéfices.}
Homogénéité, accessibilité, faible poids. \textbf{Limites.} Verbosité des classes.
\textbf{Alternatives.} CSS Modules, styled-components, bibliothèque shadcn/ui.

\subsection{Framer Motion et React Flow}
\textbf{Pourquoi.} Animations fluides et graphes interactifs pour les
fonctionnalités phares. \textbf{Comment.} Framer Motion anime l'AI Execution
Studio (pipeline, avatar, jauges) ; React Flow dessine l'Architecture Explorer.
\textbf{Bénéfices.} 60~images/s via animations sur \texttt{transform}/\texttt{opacity},
respect de \emph{prefers-reduced-motion}. \textbf{Limites.} Poids des
bibliothèques, atténué par le chargement différé. \textbf{Alternatives.} GSAP,
react-spring ; D3, Cytoscape.

\section{DevOps}
\textbf{Docker et Docker Compose.} L'ensemble est conteneurisé en huit services
(PostgreSQL, Redis, Qdrant, \emph{backend}, \emph{worker} Celery, \emph{frontend},
Prometheus, Grafana), garantissant la reproductibilité de l'environnement.
\textbf{Prometheus et Grafana} assurent la collecte et la visualisation des
métriques système~\cite{docker}.

\section{Détails d'implémentation}

\subsection{Implémentation backend}
Les routes délèguent aux services ; les modèles SQLAlchemy déclarent les onze
tables. Le listing~\ref{lst:model} illustre un modèle représentatif.

\begin{lstlisting}[style=code,language=Python,caption={Modèle SQLAlchemy (extrait)},label={lst:model}]
class AgentRun(UUIDPrimaryKeyMixin, Base):
    __tablename__ = "agent_runs"
    message_id: Mapped[UUID] = mapped_column(ForeignKey("messages.id"))
    agent_type: Mapped[str]
    status: Mapped[str]
    latency_ms: Mapped[int | None]
    input_preview: Mapped[dict] = mapped_column(JSONB)
    output_preview: Mapped[dict | None] = mapped_column(JSONB)
    error: Mapped[str | None]
\end{lstlisting}

\subsection{Implémentation frontend}
L'état serveur est géré par React Query. Le listing~\ref{lst:rq} illustre un
\emph{hook} d'inspection de trace orchestrant des requêtes parallèles.

\begin{lstlisting}[style=code,language=JavaScript,caption={Hook React Query (extrait)},label={lst:rq}]
export function useTraceInspector(messageId, includeContent, enabled) {
  const trace = useQuery({
    queryKey: ["trace", messageId, includeContent],
    queryFn: () => getRagTrace(messageId, includeContent),
    enabled,
  });
  // requetes secondaires en parallele, degradation gracieuse
  return { trace, /* retrieval, reranking, evaluation */ };
}
\end{lstlisting}

\subsection{Implémentation API}
L'API expose les domaines authentification, espaces, documents, chat (REST + SSE),
système, administration, RAGOps et traces, sous \texttt{/api/v1}. Les
\emph{endpoints} privilégiés (administration, RAGOps, traces) sont protégés par
des dépendances de rôle.

\subsection{Implémentation base de données}
Le schéma est défini en SQLAlchemy et son évolution gérée par Alembic (schéma
initial, évaluation des réponses, RBAC d'administration). Les fragments
référencent leur point vectoriel Qdrant via \texttt{vector\_id}.

\subsection{Implémentation RAG}
Le flux agentique instancie les sept agents et journalise chaque exécution. Le
listing~\ref{lst:wf} esquisse la boucle de génération en \emph{streaming}.

\begin{lstlisting}[style=code,language=Python,caption={Génération en streaming (extrait)},label={lst:wf}]
stream_generate = getattr(provider, "stream_generate", None)
if callable(stream_generate):
    async for token in stream_generate(prompt=prompt,
                                       system_prompt=SYSTEM_PROMPT):
        if token:
            raw_parts.append(token)
            yield ("token", {"text": token})
\end{lstlisting}

\subsection{Implémentation RAGOps}
Les services \texttt{ragops\_service} et \texttt{rag\_traces\_service} agrègent la
santé Qdrant, la file documentaire, la couverture d'\emph{embedding} et les
indicateurs de qualité, exposés via \texttt{/admin/ragops/*} et
\texttt{/admin/rag-traces/quality/summary}, puis restitués par la tour de contrôle.

\subsection{Implémentation traçabilité}
À la fin de chaque requête, le message et sa lignée (exécutions d'agents, passages
récupérés, évaluation, usage) sont persistés. L'inspecteur les recharge via des
requêtes parallèles et reconstruit les neuf sections d'analyse, en masquant les
secrets dans la vue de débogage.

\section{Conclusion}
Chaque brique a été retenue à l'issue d'une analyse de compromis explicite, puis
intégrée dans une implémentation cohérente : \emph{backend} en couches,
\emph{frontend} typé et modulaire, pipeline RAG agentique traçable, et
sous-systèmes d'exploitation et de visualisation. Le chapitre suivant valide ces
choix par les tests et les résultats.
